\documentclass[11pt,a4paper]{article}
\usepackage[utf8]{inputenc}
\usepackage[T1]{fontenc}
\usepackage[catalan]{babel}
\usepackage{lmodern}
\linespread{1.22}
\usepackage[a4paper,top=2.4cm,bottom=2.2cm,left=1.5cm,right=1.5cm,headheight=1.1cm,headsep=0.5cm,footskip=1cm]{geometry}
\usepackage{amsmath,amssymb,cancel}
\usepackage{icomma}
\usepackage{graphicx}
\usepackage[table]{xcolor}
\usepackage{tikz}\usetikzlibrary{calc}
\usepackage{pgfplots}\pgfplotsset{compat=1.17}
\usepackage{fontawesome5}
\usepackage{ccicons}
\usepackage{fancyhdr}
\usepackage{multicol}
\usepackage{enumitem}
\usepackage{booktabs}
\usepackage{tabularx}
\usepackage{needspace}
\usepackage[most]{tcolorbox}
\usepackage{titlesec}
\usepackage[hidelinks]{hyperref}
\hypersetup{pdftitle={Apunts de teoria. Unitat 1: Matrius i determinants. 2n de Batxillerat MII},pdfsubject={Matemàtiques II},pdfkeywords={matrius, operacions amb matrius, determinants, regla de Sarrus, adjunts, rang, matriu inversa, equacions matricials}}

% colors de la web (FentMates, bloc d'Àlgebra: blau; es manté el nom fmgreen dels apunts de les altres unitats)
\definecolor{fmnavy}{HTML}{1B2748}
\definecolor{fmgreen}{HTML}{3D6AA0}
\definecolor{fmgreenl}{HTML}{4E81BD}
\definecolor{fmyellow}{HTML}{FFC94D}
\definecolor{fmgold}{HTML}{E9A919}
\definecolor{fmaccent}{HTML}{FF7A50}
\definecolor{fmteal}{HTML}{2BB3A3}
\definecolor{fmdef}{HTML}{1D4F73}\definecolor{fmdefbg}{HTML}{EEF4F9}
\definecolor{fmprop}{HTML}{3B4252}\definecolor{fmpropbg}{HTML}{F3F4F7}
\definecolor{fmex}{HTML}{A3203A}\definecolor{fmexbg}{HTML}{FBF0F2}
\definecolor{fmmuted}{HTML}{5D6685}\definecolor{fmball0}{HTML}{E2483D}\definecolor{fmball1}{HTML}{2F6FD6}
\definecolor{fmsurface}{HTML}{EEF1F8}

\newcommand{\lt}{<}\newcommand{\gt}{>}
% fórmula en línia que no es parteix mai; si no cap en una línia, es posa centrada
\newcommand{\fmm}[1]{\setbox0=\hbox{$#1$}\ifdim\wd0>\linewidth\fmdisp{#1}\else\mbox{\box0}\fi}
% fórmula centrada; si és més ampla que la línia, s'encongeix fins que hi cap
\newcommand{\fmdisp}[1]{\[\setbox2=\hbox{$\displaystyle #1$}\ifdim\wd2>\linewidth\resizebox{\linewidth}{!}{\box2}\else\box2\fi\]}
\DeclareUnicodeCharacter{2212}{-}
\DeclareUnicodeCharacter{2225}{\ensuremath{\parallel}}\DeclareUnicodeCharacter{22A5}{\ensuremath{\perp}}\DeclareUnicodeCharacter{2226}{\ensuremath{\nparallel}}
\DeclareUnicodeCharacter{207A}{\ensuremath{^{+}}}\DeclareUnicodeCharacter{207B}{\ensuremath{^{-}}}
\setlength{\parindent}{0pt}
\setlength{\parskip}{12pt}
\setlength{\lineskip}{5pt}\setlength{\lineskiplimit}{5pt}  % més aire al voltant de les fraccions en línia
\setlength{\emergencystretch}{3em}
\setlist{itemsep=4pt,topsep=4pt}
\renewcommand{\thesection}{1.\arabic{section}}
\titleformat{\section}{\color{fmnavy}\Large\bfseries}{\colorbox{fmdef}{\color{white}\strut\ \thesection\ }}{0.6em}{}[{\color{fmdef}\titlerule[1.2pt]}]
\titlespacing*{\section}{0pt}{16pt}{8pt}
\titleformat*{\subsection}{\color{fmgreen}\large\bfseries}
\allowdisplaybreaks

% dibuix de la targeta de la unitat a la web (illustrations.js), en blanc
\newcommand{\fmcard}{\begin{tikzpicture}[x=0.062cm,y=-0.062cm,white,line width=3pt,line cap=round]
  \begin{scope}[transparency group,opacity=0.15]
  \draw (70,16)--(58,16)--(58,104)--(70,104) (170,16)--(182,16)--(182,104)--(170,104);
  \draw[line width=2pt,opacity=0.35] (78,26)--(162,92);
  \foreach \x/\y/\t in {88/40/2,120/40/$-1$,152/40/0,88/70/3,120/70/4,152/70/1,88/100/0,120/100/5,152/100/2}
    \node[anchor=base,inner sep=0pt,font=\fontsize{30}{30}\sffamily\bfseries] at (\x,\y) {\t};
  \node[anchor=base west,inner sep=0pt,font=\fontsize{26}{26}\selectfont] at (192,40) {$A^{-1}$};
  \node[anchor=base west,inner sep=0pt,font=\fontsize{26}{26}\selectfont] at (10,34) {$|A|$};
  \end{scope}
\end{tikzpicture}}
\newcommand{\fmbrand}{\raisebox{-0.25\height}{\includegraphics[height=0.5cm]{logo-fentmates.pdf}}}

% capçalera i peu
\pagestyle{fancy}\fancyhf{}
\fancyhead[L]{\small\fmbrand}
\fancyhead[R]{\small\color{fmmuted}2n Batxillerat MII · Unitat 1: Matrius i determinants}
\fancyfoot[L]{\scriptsize\color{fmmuted}\href{https://creativecommons.org/licenses/by-nc-sa/4.0/deed.ca}{\ccbyncsa}\ \ FentMates · fentmates.org · CC BY-NC-SA 4.0}
\fancyfoot[R]{\small\color{fmmuted}\thepage}
\renewcommand{\headrulewidth}{0.4pt}\renewcommand{\headrule}{\color{fmmuted!40}\hrule height \headrulewidth}
\fancypagestyle{primera}{\fancyhf{}\renewcommand{\headrulewidth}{0pt}
  \fancyfoot[L]{\scriptsize\color{fmmuted}\href{https://creativecommons.org/licenses/by-nc-sa/4.0/deed.ca}{\ccbyncsa}\ \ FentMates · fentmates.org · CC BY-NC-SA 4.0}
  \fancyfoot[R]{\small\color{fmmuted}\thepage}}

% caixes (mateixos colors que la web)
% quan una caixa es parteix entre dues pàgines, la vora del tall és discontínua
\tcbset{fmbase/.style={enhanced,breakable,arc=3mm,boxrule=1pt,left=3.5mm,right=3.5mm,top=2mm,bottom=2.5mm,before skip=10pt,after skip=10pt,fonttitle=\bfseries,
  frame style={fill=tcbcolframe},
  overlay first={\fill[tcbcolback] ([xshift=1pt]frame.south west) rectangle ([xshift=-1pt,yshift=1.2pt]frame.south east);
    \draw[tcbcolframe,line width=1pt,dash pattern=on 5pt off 3pt] (frame.south west)--(frame.south east);},
  overlay middle={\fill[tcbcolback] ([xshift=1pt]frame.south west) rectangle ([xshift=-1pt,yshift=1.2pt]frame.south east);
    \fill[tcbcolback] ([xshift=1pt]frame.north west) rectangle ([xshift=-1pt,yshift=-1.2pt]frame.north east);
    \draw[tcbcolframe,line width=1pt,dash pattern=on 5pt off 3pt] (frame.south west)--(frame.south east) (frame.north west)--(frame.north east);},
  overlay last={\fill[tcbcolback] ([xshift=1pt]frame.north west) rectangle ([xshift=-1pt,yshift=-1.2pt]frame.north east);
    \draw[tcbcolframe,line width=1pt,dash pattern=on 5pt off 3pt] (frame.north west)--(frame.north east);}}}
\newtcolorbox{fmdef}[1]{fmbase,colback=fmdefbg,colframe=fmdef,coltitle=white,title={#1}}
\newtcolorbox{fmprop}[1]{fmbase,colback=fmpropbg,colframe=fmprop,coltitle=white,title={#1},halign title=center}
\newtcolorbox{fmex}[1]{fmbase,before upper=\raggedright,colback=fmexbg,colframe=fmex,colbacktitle=fmexbg,coltitle=fmex,title={\faPencil*\ #1},
  titlerule=0.4pt,titlerule style={fmex!40,dashed}}
\newtcolorbox{fmwarn}[1]{fmbase,colback=fmsf!6,colframe=fmsf,coltitle=white,title={#1}}
\newtcolorbox{fmnote}{enhanced,breakable,colback=white,colframe=white,borderline west={3pt}{0pt}{fmaccent},arc=0pt,left=3mm,right=2mm,top=1mm,bottom=1mm,before skip=6pt,after skip=6pt,before upper=\raggedright}
\newcommand{\fmkey}[1]{\colorbox{fmyellow!45}{\strut #1}}
\newcommand{\fmuline}[1]{\textbf{\color{fmaccent}#1}}
\definecolor{fmok}{HTML}{15803D}\definecolor{fmev}{HTML}{2563EB}\definecolor{fmsf}{HTML}{C2410C}\definecolor{fmsi}{HTML}{7C3AED}\definecolor{fmas}{HTML}{A3203A}
\newcommand{\fmbadge}[2]{\tcbox[on line,colback=#1,colframe=#1,arc=2mm,boxsep=0pt,left=4pt,right=4pt,top=1.5pt,bottom=1.5pt]{\color{white}\footnotesize\bfseries #2}}
\newcommand{\fmchip}[1]{\tcbox[on line,colback=fmsf!10,colframe=fmsf!50,arc=2mm,boxsep=1pt,left=4pt,right=4pt]{#1}}
\newcommand{\fmlevel}[2]{\tcbox[on line,colback=white,colframe=fmmuted!40,arc=2mm,boxsep=2pt,halign=center]{\begin{tabular}{c}{\scriptsize\color{fmmuted}\MakeUppercase{#1}}\\#2\end{tabular}}}
% gràfiques
\pgfplotsset{fmaxis/.style={axis lines=middle,grid=major,grid style={fmmuted!15},axis line style={fmmuted,-latex},
  tick label style={font=\scriptsize,color=fmmuted},clip=true,unbounded coords=jump,restrict y to domain=-200:200,
  every axis plot/.append style={line join=round}},
  fmcurve/.style={fmaccent,line width=1.6pt,smooth}}
\tikzset{fmasy/.style={fmmuted,dashed,line width=0.9pt},fmguide/.style={fmaccent,dotted,line width=0.8pt},
  fmdot/.style={circle,fill=fmaccent,inner sep=0pt,minimum size=5pt},fmopen/.style={circle,draw=fmaccent,fill=white,line width=1pt,inner sep=0pt,minimum size=5pt}}
\newcommand{\fmpau}[1]{\tcbox[on line,colback=fmnavy,colframe=fmnavy,arc=2mm,boxsep=0pt,left=4pt,right=4pt,top=1.5pt,bottom=1.5pt]{\color{white}\scriptsize\bfseries\MakeUppercase{#1}}\ }

\begin{document}
\thispagestyle{primera}
\begin{tikzpicture}[remember picture,overlay]
  \fill[fmgreen] (current page.north west) rectangle ([yshift=-10.1cm]current page.north east);
  \begin{scope}
    \clip (current page.north west) rectangle ([yshift=-10.1cm]current page.north east);
    \node[anchor=north west,inner sep=0pt] at ([xshift=6.6cm,yshift=-0.85cm]current page.north west) {\fmcard};
  \end{scope}
  \node[anchor=north west,inner sep=0pt] at ([xshift=1.5cm,yshift=-1.1cm]current page.north west) {\tcbox[colback=white,colframe=white,arc=2.5mm,boxsep=0pt,left=7pt,right=7pt,top=5pt,bottom=5pt]{\includegraphics[height=0.75cm]{logo-fentmates.pdf}}};
\end{tikzpicture}
\vspace*{0.55cm}
\tcbox[colback=white!15!fmgreen,colframe=white!15!fmgreen,arc=2.5mm,boxsep=0pt,left=6pt,right=6pt,top=3pt,bottom=3pt]{\color{white}\small\bfseries Apunts de teoria · 2n Batxillerat · Matemàtiques II}

\vspace{4pt}
{\raggedright\color{white}\fontsize{25}{30}\selectfont\bfseries{\color{fmyellow}Unitat 1} · Matrius i determinants\par}
\vspace{6pt}
\begin{minipage}{\linewidth}\color{white!88!fmgreen}\small Teoria i exemples resolts pas a pas per entendre (i no només memoritzar) les matrius, els determinants, el rang i la matriu inversa.\end{minipage}\par
% l'índex comença sempre a la mateixa altura, per davall de la franja (com en les unitats 5-10)
\vspace*{\dimexpr 8.6cm-\pagetotal\relax}

\begin{tcolorbox}[colback=fmsurface,colframe=fmsurface,arc=3mm,left=4mm,right=4mm,top=2mm,bottom=2mm]
\small\textbf{Índex}\par\smallskip
\begin{minipage}[t]{0.49\linewidth}\small \mbox{\hyperlink{sec1}{\textbf{\color{fmgreen}1.1}\ \ Matrius. Definicions}}\\[3pt]
\mbox{\hyperlink{sec2}{\textbf{\color{fmgreen}1.2}\ \ Operacions amb matrius}}\\[3pt]
\mbox{\hyperlink{sec3}{\textbf{\color{fmgreen}1.3}\ \ Determinants}}\\[3pt]
\mbox{\hyperlink{sec4}{\textbf{\color{fmgreen}1.4}\ \ Propietats dels determinants}}\\[3pt]\end{minipage}\hfill\begin{minipage}[t]{0.49\linewidth}\small \mbox{\hyperlink{sec5}{\textbf{\color{fmgreen}1.5}\ \ Rang d'una matriu}}\\[3pt]
\mbox{\hyperlink{sec6}{\textbf{\color{fmgreen}1.6}\ \ Matriu inversa}}\\[3pt]
\mbox{\hyperlink{sec7}{\textbf{\color{fmgreen}1.7}\ \ Equacions i sistemes matricials}}\\[3pt]\end{minipage}
\end{tcolorbox}

\Needspace{18\baselineskip}
\section{Matrius. Definicions}\hypertarget{sec1}{}

\Needspace{7\baselineskip}
\begin{fmdef}{Definició: matriu d'ordre $n\times m$}
Una \textbf{matriu d'ordre o dimensió} \fmm{n\times m} és un conjunt de \fmm{n\cdot m} elements distribuïts en \fmm{n} files i \fmm{m} columnes i la representem per \fmm{A=(a_{ij})}. En general l'escriurem així:

\fmdisp{A=\begin{pmatrix} a_{11} & a_{12} & a_{13} & \cdots & a_{1m}\\ a_{21} & a_{22} & a_{23} & \cdots & a_{2m}\\ a_{31} & a_{32} & a_{33} & \cdots & a_{3m}\\ \cdots & \cdots & \cdots & \cdots & \vdots\\ a_{n1} & a_{n2} & a_{n3} & \cdots & a_{nm}\end{pmatrix}}

Cada element de la matriu l'escriurem amb dos subíndexs, \fmm{a_{ij}}, on el primer subíndex (\fmm{i}) ens indica la \textbf{fila}, i el segon subíndex (\fmm{j}), la \textbf{columna}. En la pràctica, cada element de la matriu serà un nombre real, \fmm{a_{23}=-1} (per exemple).\newline Com podem veure, una matriu \fmuline{és una taula numèrica rectangular}.
\end{fmdef}

\Needspace{7\baselineskip}
\begin{fmdef}{Definicions}
\begin{itemize}
\item Direm que una matriu és \textbf{quadrada} si \fmm{m=n}.
\item Direm que dues \textbf{matrius són iguals} quan són de la mateixa dimensió i, a més, coincideixen terme a terme:

\fmdisp{\left.\begin{array}{l} A=(a_{ij})_{m,n}\\ B=(b_{ij})_{m,n}\end{array}\right\}\; A=B \iff a_{ij}=b_{ij}}

\item Anomenem \textbf{transposada} d'una matriu \fmm{A=(a_{ij})_{m,n}} una altra matriu \fmm{A^t=(a_{ij})_{n,m}}, que s'obté en canviar en la matriu \fmm{A} les files per les columnes i les columnes per les files.

\fmdisp{A=\begin{pmatrix}1 & 2\\ 3 & 4\end{pmatrix}\qquad A^t=\begin{pmatrix}1 & 3\\ 2 & 4\end{pmatrix}}

\item Direm que una matriu \fmm{A} és \textbf{simètrica} si, a més de ser quadrada, es compleix que \fmm{A=A^t}.
\item Una matriu quadrada direm que és \textbf{triangular}, si tots els elements per davall o per dalt de la diagonal principal són iguals a zero, com la de l'exemple:

\fmdisp{\underbrace{A=\begin{pmatrix}4 & \color[HTML]{2bb3a3}0 & \color[HTML]{2bb3a3}0\\ -1 & 5 & \color[HTML]{2bb3a3}0\\ 6 & -1 & 1\end{pmatrix}}_{\text{Triangular inferior}}\quad,\quad \underbrace{B=\begin{pmatrix}3 & 3 & -6\\ \color[HTML]{2bb3a3}0 & 2 & -1\\ \color[HTML]{2bb3a3}0 & \color[HTML]{2bb3a3}0 & 9\end{pmatrix}}_{\text{Triangular superior}}}

\item Anomenem \textbf{matriu nu\lgem{}a} la matriu on tots els elements són zeros. És a dir, \fmm{a_{ij}=0\ \ \forall i,j}.
\item \textbf{Matriu diagonal} és una matriu quadrada on tots els elements que no pertanyen a la diagonal principal són zero. És a dir, \fmm{a_{ij}=0\ \ \forall i\neq j}.
\item \textbf{Matriu identitat o matriu unitat} és la matriu diagonal on tots els elements de la diagonal principal són 1, és a dir, \fmm{a_{ij}=1\ \ \forall i=j} i \fmm{a_{ij}=0\ \ \forall i\neq j}.
\end{itemize}
\end{fmdef}

\Needspace{7\baselineskip}
\begin{fmex}{Exemples de matrius}
\begin{enumerate}[label=\textbf{\arabic*.},leftmargin=*]
\item Siga la matriu \fmm{A=\begin{pmatrix}7 & 1 & 4 & 2\\ 0 & 5 & -1 & 3\\ 6 & 2 & 0 & 5\end{pmatrix}} la seua transposada serà \fmm{A^t=\begin{pmatrix}7 & 0 & 6\\ 1 & 5 & 2\\ 4 & -1 & 0\\ 2 & 3 & 5\end{pmatrix}}.
\item La matriu \fmm{B=\begin{pmatrix}1 & 6 & -5\\ 6 & 0 & 4\\ -5 & 4 & 6\end{pmatrix}} és simètrica ja que \fmm{B=B^t}.
\item \fmm{A=\begin{pmatrix}1 & 3\\ 0 & 1\end{pmatrix}} matriu triangular superior. \fmm{B=\begin{pmatrix}1 & 0 & 0\\ 2 & 4 & 0\\ 3 & -3 & -2\end{pmatrix}} matriu triangular inferior.
\item \fmm{C=\begin{pmatrix}1 & 0 & 0\\ 0 & 2 & 0\\ 0 & 0 & -3\end{pmatrix}} és una matriu diagonal.
\item \fmm{I_2=\begin{pmatrix}1 & 0\\ 0 & 1\end{pmatrix}} matriu identitat d'ordre 2. \fmm{I_3=\begin{pmatrix}1 & 0 & 0\\ 0 & 1 & 0\\ 0 & 0 & 1\end{pmatrix}} matriu identitat d'ordre 3.
\item En el seguiment de la pandèmia de la COVID-19 s'han comptabilitzat els contagis, les noves hospitalitzacions i les altes hospitalàries de quatre localitats durant una setmana.

\begin{center}\renewcommand{\arraystretch}{1.5}\small\begin{tabularx}{\linewidth}{>{\raggedright\arraybackslash}p{3.2cm}>{\raggedright\arraybackslash}X>{\raggedright\arraybackslash}X>{\raggedright\arraybackslash}X}\toprule
\textbf{} & \textbf{Contagis} & \textbf{Noves hospitalitzacions} & \textbf{Altes hospitalàries}\\\midrule
\textbf{L$_1$} & 42 & 23 & 30\\[3pt]
\textbf{L$_2$} & 35 & 15 & 16\\[3pt]
\textbf{L$_3$} & 56 & 18 & 22\\[3pt]
\textbf{L$_4$} & 12 & 3 & 5\\[3pt]
\bottomrule\end{tabularx}\end{center}

          És clar que les dades configuren una matriu 4x3:

\fmdisp{\begin{pmatrix}42 & 23 & 30\\ 35 & 15 & 16\\ 56 & 18 & 22\\ 12 & 3 & 5\end{pmatrix}}

\end{enumerate}
\end{fmex}

\Needspace{18\baselineskip}
\section{Operacions amb matrius}\hypertarget{sec2}{}

\Needspace{7\baselineskip}
\begin{fmdef}{Suma/resta de matrius}
Per a poder \textbf{sumar/restar} dues matrius cal que aquestes siguen del mateix ordre, i sols haurem de sumar/restar els elements que ocupen la mateixa posició.

\fmdisp{A\pm B=(a_{ij})\pm(b_{ij})=(a_{ij}\pm b_{ij});\ \forall i,j}

\end{fmdef}

\Needspace{7\baselineskip}
\begin{fmprop}{Propietats de la suma}
\begin{enumerate}[label=\alph*)]
\item Associativa: \fmm{A+(B+C)=(A+B)+C}
\item Commutativa: \fmm{A+B=B+A}
\item Element neutre: \fmm{A+0=A=0+A}
\item Element simètric: \fmm{A+(-A)=0=(-A)+A}. La matriu \fmm{(-A)} s'anomena \textbf{matriu oposada} d'\fmm{A} i està formada per tots els elements oposats de la matriu \fmm{A}, és a dir, els elements d'\fmm{A} amb signe contrari.
\end{enumerate}
\end{fmprop}

\Needspace{7\baselineskip}
\begin{fmex}{Exemples de suma/resta de matrius}
Siguen les matrius \fmm{A=\begin{pmatrix}5 & 3 & 1\\ 6 & 4 & 1\end{pmatrix}} i \fmm{B=\begin{pmatrix}1 & 2 & 1\\ 5 & 3 & 0\end{pmatrix}}, calculem:

\fmdisp{A+B=\begin{pmatrix}5 & 3 & 1\\ 6 & 4 & 1\end{pmatrix}+\begin{pmatrix}1 & 2 & 1\\ 5 & 3 & 0\end{pmatrix}=\begin{pmatrix}5+1 & 3+2 & 1+1\\ 6+5 & 4+3 & 1+0\end{pmatrix}=\begin{pmatrix}6 & 5 & 2\\ 11 & 7 & 1\end{pmatrix}}

\fmdisp{A-B=\begin{pmatrix}5 & 3 & 1\\ 6 & 4 & 1\end{pmatrix}-\begin{pmatrix}1 & 2 & 1\\ 5 & 3 & 0\end{pmatrix}=\begin{pmatrix}5-1 & 3-2 & 1-1\\ 6-5 & 4-3 & 1-0\end{pmatrix}=\begin{pmatrix}4 & 1 & 0\\ 1 & 1 & 1\end{pmatrix}}

Les matrius \fmm{\begin{pmatrix}2 & 1 & 7 & 4\\ -3 & 2 & -5 & 6\end{pmatrix}} i \fmm{\begin{pmatrix}4 & 7 & 0\\ -4 & 1 & -2\end{pmatrix}} no poden sumar-se perquè no tenen el mateix ordre.
\end{fmex}

\Needspace{7\baselineskip}
\begin{fmdef}{Producte d'una matriu per un nombre real}
Donada una matriu \fmm{A} i un nombre real \fmm{k}, el \textbf{producte} \fmm{k\cdot A} és una altra matriu del mateix ordre, els elements de la qual són el producte del nombre \fmm{k} pels elements de la matriu \fmm{A}:

\fmdisp{k\cdot A=k\cdot(a_{ij})=(k\cdot a_{ij})\ \ \forall i,j}

\end{fmdef}

\Needspace{7\baselineskip}
\begin{fmprop}{Propietats}
Siguen \fmm{k} i \fmm{n} dos nombres reals i \fmm{A} i \fmm{B} dues matrius del mateix ordre:

\begin{enumerate}[label=\alph*)]
\item Associativa: \fmm{k\cdot(n\cdot A)=(k\cdot n)\cdot A}
\item Distributiva respecte de la suma de matrius: \fmm{k\cdot(A+B)=k\cdot A+k\cdot B}
\item Distributiva respecte de la suma de nombres: \fmm{(n+k)\cdot A=n\cdot A+k\cdot A}
\item Element neutre: \fmm{1\cdot A=A}
\end{enumerate}
\end{fmprop}

\Needspace{7\baselineskip}
\begin{fmex}{Exemple}
Donada la matriu \fmm{A=\begin{pmatrix}1 & 2\\ 3 & 0\\ 1 & -1\end{pmatrix}}, calculem:

\fmdisp{2\cdot A=2\cdot\begin{pmatrix}1 & 2\\ 3 & 0\\ 1 & -1\end{pmatrix}=\begin{pmatrix}2 & 4\\ 6 & 0\\ 2 & -2\end{pmatrix}}

\end{fmex}

\Needspace{7\baselineskip}
\begin{fmdef}{Producte de matrius}
Perquè dues matrius \fmm{A} (d'ordre \fmm{m\times n}) i \fmm{B} (d'ordre \fmm{n\times p}) puguen multiplicar-se, \fmm{A\cdot B}, cal que \fmkey{el nombre de columnes de la primera coincidisca amb el nombre de files de la segona}.\newline En tal cas, el producte \fmm{A\cdot B=C} és una altra matriu els elements de la qual s'obtenen multiplicant cada vector fila de la primera matriu per cada vector columna de la segona matriu, de la manera següent:

\fmdisp{C=A\cdot B=(a_{ij})\cdot(b_{ij})=c_{ij}=\sum_{k=1}^{n}a_{ik}\cdot b_{kj}}

\end{fmdef}

\Needspace{7\baselineskip}
\begin{fmprop}{Propietats del producte}
\begin{enumerate}[label=\alph*)]
\item Associativa: \fmm{A\cdot(B\cdot C)=(A\cdot B)\cdot C}
\item Element neutre: \fmm{A\cdot I=I\cdot A=A}
\item Distributiva respecte de la suma: \fmm{A\cdot(B+C)=A\cdot B+A\cdot C}
\item Transposada del producte: \fmm{(A\cdot B)^t=B^t\cdot A^t}
\end{enumerate}
\end{fmprop}

\Needspace{7\baselineskip}
\begin{fmex}{Exemple}
Donades les matrius \fmm{A=\begin{pmatrix}5 & 3 & 1\\ 6 & 4 & 1\end{pmatrix}} i \fmm{B=\begin{pmatrix}1 & 2\\ 3 & 0\\ 1 & -1\end{pmatrix}}, calculem \fmm{A\cdot B} i \fmm{B\cdot A}:

\fmdisp{A\cdot B=\begin{pmatrix}5 & 3 & 1\\ 6 & 4 & 1\end{pmatrix}\cdot\begin{pmatrix}1 & 2\\ 3 & 0\\ 1 & -1\end{pmatrix}=\begin{pmatrix}5\cdot1+3\cdot3+1\cdot1 & 5\cdot2+3\cdot0+1\cdot(-1)\\ 6\cdot1+4\cdot3+1\cdot1 & 6\cdot2+4\cdot0+1\cdot(-1)\end{pmatrix}=\begin{pmatrix}15 & 9\\ 19 & 11\end{pmatrix}}

\fmdisp{B\cdot A=\begin{pmatrix}1 & 2\\ 3 & 0\\ 1 & -1\end{pmatrix}\cdot\begin{pmatrix}5 & 3 & 1\\ 6 & 4 & 1\end{pmatrix}=\begin{pmatrix}1\cdot5+2\cdot6 & 1\cdot3+2\cdot4 & 1\cdot1+2\cdot1\\ 3\cdot5+0\cdot6 & 3\cdot3+0\cdot4 & 3\cdot1+0\cdot1\\ 1\cdot5+(-1)\cdot6 & 1\cdot3+(-1)\cdot4 & 1\cdot1+(-1)\cdot1\end{pmatrix}=\begin{pmatrix}17 & 11 & 3\\ 15 & 9 & 3\\ -1 & -1 & 0\end{pmatrix}}

\end{fmex}

\Needspace{7\baselineskip}
\begin{fmprop}{Compte amb les no propietats}
Com podem observar, no hem parlat de la propietat commutativa en el producte de matrius i el motiu és que no es compleix sempre. En l'exemple anterior podem comprovar que no ens ha donat el mateix resultat \fmm{A\cdot B} que \fmm{B\cdot A}, és a dir:

\fmdisp{A\cdot B\neq B\cdot A}

Açò fa que altres propietats que utilitzem amb nombres no es complisquen tampoc per a matrius:

\begin{itemize}
\item \fmm{(A+B)\cdot(A-B)\neq A^2-B^2}
\item \fmm{(A\pm B)^2\neq A^2\pm 2AB+B^2}
\item \fmm{A\cdot B=0 \not\Rightarrow A=0\ \text{o}\ B=0}
\item \fmm{A\cdot B=A\cdot C \not\Rightarrow B=C}
\end{itemize}

\textbf{En el cas que \fmm{A\cdot B=B\cdot A} diem que \fmm{A} i \fmm{B} són permutables.}

\vspace*{0.2cm}

\begin{fmnote}\textbf{NOTA:} Cal destacar el cas en què una de les dues matrius siga la matriu identitat, ja que aquesta és permutable amb qualsevol matriu (per la definició de matriu identitat):

\fmdisp{A\cdot I=I\cdot A}

Per tant, en aquest cas es compliran les identitats notables:

\begin{itemize}
\item \fmm{(A+I)\cdot(A-I)=A^2-I^2=A^2-I}, ja que:
\vspace*{-0.2cm}
\fmdisp{(A+I)\cdot(A-I)=A\cdot A-A\cdot I+I\cdot A-I\cdot I\overset{A\cdot I=I\cdot A}{=}A^2-A\cdot I+A\cdot I-I^2=A^2-I^2=A^2-I}

\item \fmm{(A\pm I)^2=A^2\pm2AI+I^2=A^2\pm2A+I} ja que:
\vspace*{-0.2cm}
\fmdisp{(A\pm I)^2=(A\pm I)\cdot(A\pm I)=A\cdot A\pm A\cdot I\pm I\cdot A+I\cdot I\overset{A\cdot I=I\cdot A}{=}A^2\pm A\cdot I\pm A\cdot I+I^2=}
\fmdisp{=A^2\pm2AI+I=A^2\pm2A+I}

\end{itemize}\end{fmnote}
\end{fmprop}

\Needspace{7\baselineskip}
\begin{fmdef}{Potència de matrius}
Per calcular la \textbf{potència d'una matriu} quadrada d'ordre \fmm{n}, es multiplica la matriu per si mateixa tantes vegades com indique l'exponent, és a dir:

\fmdisp{A^n=A\cdot A\cdot A\cdots A\ \ (n\ \text{vegades})}

La \textbf{potència enèsima} d'una matriu \fmm{A} és la matriu que representa totes les potències d'\fmm{A} en funció del valor que prenga l'exponent. L'escriurem \fmm{A^n}.
\end{fmdef}

\Needspace{7\baselineskip}
\begin{fmex}{Exemple de potències de matrius}
Donada la matriu \fmm{B=\begin{pmatrix}1 & 1\\ 1 & 1\end{pmatrix}}, calcularem \fmm{B^n} i \fmm{B^{12}}:

\fmdisp{B^2=\begin{pmatrix}1 & 1\\ 1 & 1\end{pmatrix}\cdot\begin{pmatrix}1 & 1\\ 1 & 1\end{pmatrix}=\begin{pmatrix}2 & 2\\ 2 & 2\end{pmatrix}\qquad B^3=\begin{pmatrix}1 & 1\\ 1 & 1\end{pmatrix}\cdot\begin{pmatrix}2 & 2\\ 2 & 2\end{pmatrix}=\begin{pmatrix}4 & 4\\ 4 & 4\end{pmatrix}=\begin{pmatrix}2^2 & 2^2\\ 2^2 & 2^2\end{pmatrix}}

\fmdisp{B^4=B\cdot B^3=\begin{pmatrix}1 & 1\\ 1 & 1\end{pmatrix}\cdot\begin{pmatrix}4 & 4\\ 4 & 4\end{pmatrix}=\begin{pmatrix}8 & 8\\ 8 & 8\end{pmatrix}=\begin{pmatrix}2^3 & 2^3\\ 2^3 & 2^3\end{pmatrix}}

Aleshores podem deduir que:

\fmdisp{\boxed{B^n=\begin{pmatrix}2^{n-1} & 2^{n-1}\\ 2^{n-1} & 2^{n-1}\end{pmatrix}}\ \ \text{i després calcular}\ \ \boxed{B^{12}=\begin{pmatrix}2^{11} & 2^{11}\\ 2^{11} & 2^{11}\end{pmatrix}=\begin{pmatrix}2048 & 2048\\ 2048 & 2048\end{pmatrix}}}

\end{fmex}

\Needspace{7\baselineskip}
\begin{fmprop}{Potències de matrius relacionades amb la unitat}
Si en calcular les potències d'una matriu \fmm{A} ens apareix la matriu identitat, \fmm{I}, aleshores les potències successives seran repetició de les anteriors i es podrà determinar, sense problema, la matriu \fmm{A^n} per a qualsevol valor de \fmm{n}. Vegem un exemple.

Donada la matriu \fmm{A=\begin{pmatrix}2 & 3\\ -1 & -2\end{pmatrix}}, determinar \fmm{A^n} i \fmm{A^{51}}:

Calculem \fmm{A^2=\begin{pmatrix}2 & 3\\ -1 & -2\end{pmatrix}\cdot\begin{pmatrix}2 & 3\\ -1 & -2\end{pmatrix}=\begin{pmatrix}1 & 0\\ 0 & 1\end{pmatrix}}; vegem ara \fmm{A^3=A^2\cdot A=I\cdot A=A}

Calculem també \fmm{A^4=A^2\cdot A^2=I\cdot I=I}. Aleshores podem deduir que:

\fmdisp{\boxed{A^n=\begin{cases} I, & \text{si $n$ parell}\\[4pt] A, & \text{si $n$ imparell}\end{cases}}\quad\text{i per tant}\quad\boxed{A^{51}=A}}

\end{fmprop}

\Needspace{18\baselineskip}
\section{Determinants}\hypertarget{sec3}{}

\Needspace{7\baselineskip}
\begin{fmdef}{Determinants}
Els \textbf{determinants} \fmuline{són nombres} associats a les matrius quadrades d'ordre \fmm{n} que es calculen de diferents formes segons siga l'ordre de la matriu.\newline Donada una matriu quadrada \fmm{A} d'ordre \fmm{n}, el seu determinant l'escriurem:

\fmdisp{\boxed{\det(A)=|A|}}

\end{fmdef}

\Needspace{7\baselineskip}
\begin{fmprop}{Càlcul del determinant d'una matriu segons el seu ordre}
\begin{enumerate}
\item Si \fmm{n=1}; \fmm{A=(a_{11})} i el determinant serà: \fmm{|A|=a_{11}}.
\item Si \fmm{n=2}; \fmm{A=\begin{pmatrix}a_{11} & a_{12}\\ a_{21} & a_{22}\end{pmatrix}} i el determinant serà: \fmm{|A|=a_{11}\cdot a_{22}-a_{12}\cdot a_{21}}
\item Si \fmm{n=3}; \fmm{A=\begin{pmatrix}a_{11} & a_{12} & a_{13}\\ a_{21} & a_{22} & a_{23}\\ a_{31} & a_{32} & a_{33}\end{pmatrix}} en aquest cas utilitzarem l'anomenada \textbf{regla de Sarrus}:

\fmdisp{\begin{vmatrix}a_{11} & a_{12} & a_{13}\\ a_{21} & a_{22} & a_{23}\\ a_{31} & a_{32} & a_{33}\end{vmatrix}=a_{11}a_{22}a_{33}+a_{12}a_{23}a_{31}+a_{13}a_{21}a_{32}-a_{13}a_{22}a_{31}-a_{12}a_{21}a_{33}-a_{11}a_{23}a_{32}}

La meitat dels sumands tenen \textbf{signe +}, i l'altra meitat, \textbf{signe -}. Aquests sis sumands es recorden fàcilment amb la regla mnemotècnica següent, anomenada regla de Sarrus:

\begin{center}\begin{tikzpicture}[x=0.0205cm,y=-0.0205cm,line join=round]\draw[line width=0.9pt] (8,12)--(8,160) (202,12)--(202,160);\draw[fmok,line width=1.6pt,opacity=0.85] (35,30)--(105,85)--(175,140);\draw[fmok,line width=1.6pt,opacity=0.85] (105,30)--(175,85)--(35,140)--(105,30);\draw[fmok,line width=1.6pt,opacity=0.85] (175,30)--(35,85)--(105,140)--(175,30);\node[fill=white,fill opacity=0.8,text opacity=1,inner sep=1pt] at (35,30) {$a_{11}$};\node[fill=white,fill opacity=0.8,text opacity=1,inner sep=1pt] at (105,30) {$a_{12}$};\node[fill=white,fill opacity=0.8,text opacity=1,inner sep=1pt] at (175,30) {$a_{13}$};\node[fill=white,fill opacity=0.8,text opacity=1,inner sep=1pt] at (35,85) {$a_{21}$};\node[fill=white,fill opacity=0.8,text opacity=1,inner sep=1pt] at (105,85) {$a_{22}$};\node[fill=white,fill opacity=0.8,text opacity=1,inner sep=1pt] at (175,85) {$a_{23}$};\node[fill=white,fill opacity=0.8,text opacity=1,inner sep=1pt] at (35,140) {$a_{31}$};\node[fill=white,fill opacity=0.8,text opacity=1,inner sep=1pt] at (105,140) {$a_{32}$};\node[fill=white,fill opacity=0.8,text opacity=1,inner sep=1pt] at (175,140) {$a_{33}$};\node[text=fmok,font=\Large\bfseries] at (105,188) {$+$};\end{tikzpicture}\qquad\qquad\begin{tikzpicture}[x=0.0205cm,y=-0.0205cm,line join=round]\draw[line width=0.9pt] (8,12)--(8,160) (202,12)--(202,160);\draw[fmas,line width=1.6pt,opacity=0.85] (175,30)--(105,85)--(35,140);\draw[fmas,line width=1.6pt,opacity=0.85] (105,30)--(35,85)--(175,140)--(105,30);\draw[fmas,line width=1.6pt,opacity=0.85] (35,30)--(175,85)--(105,140)--(35,30);\node[fill=white,fill opacity=0.8,text opacity=1,inner sep=1pt] at (35,30) {$a_{11}$};\node[fill=white,fill opacity=0.8,text opacity=1,inner sep=1pt] at (105,30) {$a_{12}$};\node[fill=white,fill opacity=0.8,text opacity=1,inner sep=1pt] at (175,30) {$a_{13}$};\node[fill=white,fill opacity=0.8,text opacity=1,inner sep=1pt] at (35,85) {$a_{21}$};\node[fill=white,fill opacity=0.8,text opacity=1,inner sep=1pt] at (105,85) {$a_{22}$};\node[fill=white,fill opacity=0.8,text opacity=1,inner sep=1pt] at (175,85) {$a_{23}$};\node[fill=white,fill opacity=0.8,text opacity=1,inner sep=1pt] at (35,140) {$a_{31}$};\node[fill=white,fill opacity=0.8,text opacity=1,inner sep=1pt] at (105,140) {$a_{32}$};\node[fill=white,fill opacity=0.8,text opacity=1,inner sep=1pt] at (175,140) {$a_{33}$};\node[text=fmas,font=\Large\bfseries] at (105,188) {$-$};\end{tikzpicture}\end{center}
\item Si \fmm{n>3}, per calcular els determinants ho farem utilitzant els \textbf{adjunts d'una matriu}. Aquest mètode també el podem utilitzar per a matrius d'ordre menor, però els mètodes anteriors són més ràpids en la pràctica.\newline Primer veurem l'explicació teòrica i després amb un exemple comprovarem que no és tan complicat com sembla.
\end{enumerate}
\end{fmprop}

\Needspace{7\baselineskip}
\begin{fmdef}{Menor complementari i adjunt}
Donada \fmm{A} una matriu quadrada d'ordre \fmm{n}, el \textbf{menor complementari} \fmm{\alpha_{ij}} de l'element \fmm{a_{ij}} és el determinant d'ordre \fmm{n-1} que resulta d'eliminar la fila \fmm{i} i la columna \fmm{j} de la matriu \fmm{A}.\newline Anomenem \textbf{adjunt} de l'element \fmm{a_{ij}} \fmm{A_{ij}=\alpha_{ij}\cdot(-1)^{i+j}}.\newline Amb aquestes definicions ja podem dir com podem calcular el determinant de la matriu \fmm{A}:

\begin{fmnote}El valor d'un determinant és igual a la suma dels productes dels elements d'una fila o columna pels seus adjunts corresponents:\end{fmnote}

\fmdisp{\boxed{|A_n|=\sum_{k=1}^{n}a_{ik}\cdot A_{ik}=\sum_{k=1}^{n}a_{kj}\cdot A_{kj}}}

\end{fmdef}

\Needspace{7\baselineskip}
\begin{fmex}{Exemples de càlcul de determinants}
Vegem amb exemples concrets com apliquem els diferents mètodes per calcular els determinants:

\begin{enumerate}[label=\textbf{\arabic*.},leftmargin=*]
\item Siga la matriu \fmm{A=\begin{pmatrix}-1 & 2\\ 3 & 1\end{pmatrix}}, el seu determinant serà: \fmm{|A|=(-1)\cdot1-2\cdot3=-7}.
\item Siga la matriu \fmm{A=\begin{pmatrix}2 & 1 & 3\\ 4 & -1 & 2\\ -3 & 4 & 2\end{pmatrix}}; utilitzem la regla de Sarrus:

\fmdisp{|A|=\begin{vmatrix}2 & 1 & 3\\ 4 & -1 & 2\\ -3 & 4 & 2\end{vmatrix}=2\cdot(-1)\cdot2+4\cdot4\cdot3+1\cdot2\cdot(-3)-[3\cdot(-1)\cdot(-3)+2\cdot4\cdot2+1\cdot4\cdot2]=}

\fmdisp{=-4+48-6-[9+16+8]=38-33=\boxed{5}}

\item Ara utilitzarem els adjunts d'una matriu per calcular el determinant, però ho farem en una matriu de 3x3 per no fer-ho molt llarg.\newline
          Siga la matriu \fmm{A=\begin{pmatrix}1 & 2 & 3\\ 5 & 4 & 1\\ 0 & 4 & 2\end{pmatrix}}, calculem el seu determinant (desenvolupem per la primera columna):

\fmdisp{|A|=\begin{vmatrix}1 & 2 & 3\\ 5 & 4 & 1\\ 0 & 4 & 2\end{vmatrix}=1\cdot A_{11}+5\cdot A_{21}+\cancel{0\cdot A_{31}}^{\,0}}

\fmdisp{\alpha_{11}=\begin{vmatrix}4 & 1\\ 4 & 2\end{vmatrix}=8-4=4\ \rightarrow\ A_{11}=\alpha_{11}\cdot(-1)^{1+1}=4\cdot1=4}

\fmdisp{\alpha_{21}=\begin{vmatrix}2 & 3\\ 4 & 2\end{vmatrix}=4-12=-8\ \rightarrow\ A_{21}=\alpha_{21}\cdot(-1)^{2+1}=(-8)\cdot(-1)=8}

\fmdisp{|A|=1\cdot4+5\cdot8=4+40=\boxed{44}}

Com podem veure, sempre ens resultarà més senzill si desenvolupem per adjunts agafant una fila o columna que tinga zeros, ja que eixos termes s'anu\lgem{}en.
\end{enumerate}
\end{fmex}

\Needspace{18\baselineskip}
\section{Propietats dels determinants}\hypertarget{sec4}{}

\Needspace{7\baselineskip}
\begin{fmprop}{Propietats dels determinants}
\begin{enumerate}
\item Si a una línia d'un determinant se li suma o resta una combinació lineal de línies para\lgem{}eles, el valor del determinant no canvia.
\item Si en un determinant s'intercanvien les files per les columnes, el seu valor no varia: \fmm{|A|=|A^t|}
\item Si intercanviem entre si dues files para\lgem{}eles, el valor del determinant canvia de signe.
\item Si el determinant conté una fila formada per zeros, aleshores el seu valor és zero.
\item Si el determinant conté dues files (para\lgem{}eles) iguals o proporcionals, aleshores el valor del determinant és zero.
\item Si es multiplica un determinant per un nombre, aquest multiplica qualsevol de les seues línies, però sols una d'elles. I viceversa, si tots els elements d'una mateixa línia són múltiples d'un mateix nombre, el valor del determinant és múltiple d'aquest nombre.
\item El determinant del producte de dues matrius és igual al producte dels determinants de cada matriu: \fmm{|A\cdot B|=|A|\cdot|B|}.
\item El determinant de la potència d'una matriu és igual a la potència del determinant de la matriu: \fmm{|A^n|=|A|^n}.
\item Si els elements d'una línia es descomponen en dos sumands, el determinant pot descompondre's en la suma de dos determinants en què els elements romanen iguals en ambdós, amb l'única excepció de la línia que es descompon en dos sumands, els quals estaran un en cada determinant de la suma.
\end{enumerate}
\end{fmprop}

\Needspace{7\baselineskip}
\begin{fmex}{Exemples de les propietats}
\begin{enumerate}[label=\textbf{\arabic*.},leftmargin=*]
\item \textbf{Propietat 1:}

\fmdisp{\begin{vmatrix}2 & 1 & 1 & 1\\ 1 & -3 & 2 & -3\\ 1 & 2 & 0 & 1\\ 1 & 2 & -1 & 3\end{vmatrix}\underset{\substack{F_2'=F_2+2F_4\\ F_4'=F_4+F_1}}{=}\begin{vmatrix}2 & 1 & \color[HTML]{ff7a50}1 & 1\\ 3 & 1 & 0 & 3\\ 1 & 2 & 0 & 1\\ 3 & 3 & 0 & 4\end{vmatrix}=(-1)^{1+3}\cdot1\cdot\begin{vmatrix}3 & 1 & 3\\ 1 & 2 & 1\\ 3 & 3 & 4\end{vmatrix}=(-1)^{1+3}\cdot1\cdot5=\boxed{5}}

\fmdisp{\begin{vmatrix}3 & 1 & 3\\ 1 & 2 & 1\\ 3 & 3 & 4\end{vmatrix}=24+9+3-(18+9+4)=36-31=5}

\item \textbf{Propietat 6:}

\fmdisp{\begin{vmatrix}15 & 3 & 15\\ 1 & 2 & 1\\ 2 & 6 & 4\end{vmatrix}=3\cdot\begin{vmatrix}5 & 1 & 5\\ 1 & 2 & 1\\ 2 & 6 & 4\end{vmatrix}=3\cdot2\begin{vmatrix}5 & 1 & 5\\ 1 & 2 & 1\\ 1 & 3 & 2\end{vmatrix}=6\cdot[20+15+1-(10+15+2)]=\boxed{54}}

\item \textbf{Propietat 9:}

\fmdisp{\begin{vmatrix}a+1 & a+2 & a+3\\ 1 & 2 & 2\\ 4 & 3 & 4\end{vmatrix}=\begin{vmatrix}a & a & a\\ 1 & 2 & 2\\ 4 & 3 & 4\end{vmatrix}+\begin{vmatrix}1 & 2 & 3\\ 1 & 2 & 2\\ 4 & 3 & 4\end{vmatrix}=}

\fmdisp{=[8a+3a+8a-(8a+6a+4a)]+[8+9+16-(24+6+8)]=\boxed{a-5}}

\item \textbf{Propietat 8:} Siga la matriu \fmm{A=\begin{pmatrix}1 & 2 & 1\\ 0 & 3 & 2\\ 0 & 2 & 2\end{pmatrix}}, comprovem si \fmm{|A^4|=|A|^4}:

\fmdisp{|A|=\begin{vmatrix}1 & 2 & 1\\ 0 & 3 & 2\\ 0 & 2 & 2\end{vmatrix}=1\cdot(6-4)=2\ \Rightarrow\ |A|^4=2^4=\boxed{16}}

\fmdisp{A^2=\begin{pmatrix}1 & 2 & 1\\ 0 & 3 & 2\\ 0 & 2 & 2\end{pmatrix}\cdot\begin{pmatrix}1 & 2 & 1\\ 0 & 3 & 2\\ 0 & 2 & 2\end{pmatrix}=\begin{pmatrix}1 & 10 & 7\\ 0 & 13 & 10\\ 0 & 10 & 8\end{pmatrix}}

\fmdisp{A^4=A^2\cdot A^2=\begin{pmatrix}1 & 10 & 7\\ 0 & 13 & 10\\ 0 & 10 & 8\end{pmatrix}\cdot\begin{pmatrix}1 & 10 & 7\\ 0 & 13 & 10\\ 0 & 10 & 8\end{pmatrix}=\begin{pmatrix}1 & 210 & 163\\ 0 & 269 & 210\\ 0 & 210 & 164\end{pmatrix}}

Aleshores:

\fmdisp{|A^4|=\begin{vmatrix}1 & 210 & 163\\ 0 & 269 & 210\\ 0 & 210 & 164\end{vmatrix}=1\cdot(269\cdot164-210\cdot210)=\boxed{16}}

\item Siga \fmm{A=\begin{pmatrix}a & b & c\\ d & e & f\\ g & h & i\end{pmatrix}}, tal que \fmm{|A|=5}, calcular \fmm{|B|=\begin{vmatrix}-2d-2g & f+i & e+h\\ -6a & 3c & 3b\\ 2g & -i & -h\end{vmatrix}}.

Per la propietat 9 descomponem la primera fila:

\fmdisp{|B|=\begin{vmatrix}-2d-2g & f+i & e+h\\ -6a & 3c & 3b\\ 2g & -i & -h\end{vmatrix}=\begin{vmatrix}-2d & f & e\\ -6a & 3c & 3b\\ 2g & -i & -h\end{vmatrix}+\underbrace{\begin{vmatrix}-2g & i & h\\ -6a & 3c & 3b\\ 2g & -i & -h\end{vmatrix}}_{0}=30+0=\boxed{30}\ \ \text{ja que:}}

\fmdisp{\begin{vmatrix}-2d & f & e\\ -6a & 3c & 3b\\ 2g & -i & -h\end{vmatrix}\underset{(-1)\cdot C_1}{=}(-1)\cdot\begin{vmatrix}2d & f & e\\ 6a & 3c & 3b\\ -2g & -i & -h\end{vmatrix}\underset{(-1)\cdot F_3}{=}\begin{vmatrix}2d & f & e\\ 6a & 3c & 3b\\ 2g & i & h\end{vmatrix}\underset{2\cdot C_1}{=}2\cdot\begin{vmatrix}d & f & e\\ 3a & 3c & 3b\\ g & i & h\end{vmatrix}=}

\fmdisp{\underset{3\cdot F_2}{=}2\cdot3\cdot\begin{vmatrix}d & f & e\\ a & c & b\\ g & i & h\end{vmatrix}\underset{C_2\Leftrightarrow C_3}{=}-6\cdot\begin{vmatrix}d & e & f\\ a & b & c\\ g & h & i\end{vmatrix}\underset{F_1\Leftrightarrow F_2}{=}6\cdot\begin{vmatrix}a & b & c\\ d & e & f\\ g & h & i\end{vmatrix}=6\cdot5=\boxed{30}}

i com:

\fmdisp{\begin{vmatrix}-2g & i & h\\ -6a & 3c & 3b\\ 2g & -i & -h\end{vmatrix}\underset{(-1)F_3}{=}(-1)\cdot\begin{vmatrix}-2g & i & h\\ -6a & 3c & 3b\\ -2g & i & h\end{vmatrix}\underset{F_1=F_3}{=}(-1)\cdot0=0}

\end{enumerate}
\end{fmex}

\Needspace{18\baselineskip}
\section{Rang d'una matriu}\hypertarget{sec5}{}

Abans de definir el rang d'una matriu necessitem introduir alguns conceptes nous.

\Needspace{7\baselineskip}
\begin{fmdef}{Matriu singular}
Donada una matriu quadrada \fmm{A} amb determinant igual a zero, direm que és una \textbf{matriu singular}:

\fmdisp{A\ \text{singular}\iff|A|=0}

\end{fmdef}

\Needspace{7\baselineskip}
\begin{fmdef}{Matriu regular}
Donada una matriu quadrada \fmm{A} amb determinant diferent de zero, direm que és una \textbf{matriu regular}:

\fmdisp{A\ \text{regular}\iff|A|\neq0}

\end{fmdef}

\Needspace{7\baselineskip}
\begin{fmdef}{Vector fila i columna}
Un \textbf{vector fila} d'una matriu \fmm{A} és qualsevol de les seues files.\newline Un \textbf{vector columna} d'una matriu \fmm{A} és qualsevol de les seues columnes.

\fmuline{Per exemple}: donada una matriu \fmm{A=\begin{pmatrix}1 & 3\\ 2 & 6\end{pmatrix}}, direm que \fmm{(1,3)} i \fmm{(2,6)} són vectors fila de la matriu \fmm{A}. I \fmm{(1,2)} i \fmm{(3,6)} són vectors columna de la mateixa matriu.
\end{fmdef}

\Needspace{7\baselineskip}
\begin{fmdef}{Dependència i independència lineal}
\begin{itemize}
\item Es diu que un vector \fmm{\vec w} és \textbf{linealment dependent} dels vectors \fmm{\vec v_1,\vec v_2,\vec v_3,\cdots,\vec v_n} si \fmm{\vec w} és combinació lineal d'ells, és a dir:

\fmdisp{\vec w=\lambda_1\cdot\vec v_1+\lambda_2\cdot\vec v_2+\cdots+\lambda_n\cdot\vec v_n,\quad /\ \lambda_i\in\mathbb{R}\ \ \forall i\in\mathbb{N}}

\item Es diu que un vector \fmm{\vec w} és \textbf{linealment independent} dels vectors \fmm{\vec v_1,\vec v_2,\vec v_3,\cdots,\vec v_n} si \fmm{\vec w} no és combinació lineal d'ells, és a dir:

\fmdisp{\vec w\neq\lambda_1\cdot\vec v_1+\lambda_2\cdot\vec v_2+\cdots+\lambda_n\cdot\vec v_n,\quad /\ \lambda_i\in\mathbb{R}\ \ \forall i\in\mathbb{N}}

\end{itemize}
\end{fmdef}

\Needspace{7\baselineskip}
\begin{fmex}{Exemples}
\begin{enumerate}[label=\textbf{\arabic*.},leftmargin=*]
\item \textbf{És (2,0,10) linealment independent de (1,1,4) i (1,2,3)?}\newline
          Ens preguntem si existeixen nombres \fmm{\lambda} i \fmm{\mu} de manera que: \fmm{(2,0,10)=\lambda(1,1,4)+\mu(1,2,3)}

Resolem aquest sistema: \fmm{0=\lambda+2\mu\rightarrow\lambda=-2\mu}

\fmdisp{2=\lambda+\mu\ \xrightarrow[\lambda=-2\mu]{}\ 2=-2\mu+\mu=-\mu\ \rightarrow\ \boxed{\mu=-2}\ \rightarrow\ \boxed{\lambda=4}}

\fmm{10=4\lambda+3\mu\longrightarrow10=4\cdot4+3\cdot(-2)=10\ \rightarrow} el vector (2,0,10) és \textbf{l.d.} de (1,1,4) i (1,2,3).
\item \textbf{És (0,3,7) linealment independent de (1,1,4) i (1,2,3)?}\newline
          Ens preguntem si existeixen nombres \fmm{\lambda} i \fmm{\mu} de manera que: \fmm{(0,3,7)=\lambda(1,1,4)+\mu(1,2,3)}

Resolem aquest sistema: \fmm{0=\lambda+\mu\rightarrow\mu=-\lambda}

\fmdisp{3=\lambda+2\mu\ \xrightarrow[\mu=-\lambda]{}\ 3=\lambda-2\lambda=-\lambda\ \rightarrow\ \boxed{\lambda=-3}\ \rightarrow\ \boxed{\mu=3}}

\fmm{7=4\lambda+3\mu\longrightarrow4\cdot(-3)+3\cdot3=-12+9=-3\neq7\ \rightarrow} el vector (0,3,7) és \textbf{l.i.} de (1,1,4) i (1,2,3), ja que no podem trobar dos nombres \fmm{(\lambda,\mu)} que formen eixa descomposició del vector.
\end{enumerate}
\end{fmex}

\Needspace{7\baselineskip}
\begin{fmdef}{Conjunt de vectors linealment dependents i independents}
\begin{enumerate}
\item Un conjunt de vectors és \textbf{linealment dependent} si algun d'ells és dependent dels altres.
\item Un conjunt de vectors és \textbf{linealment independent} si cap d'ells és dependent dels altres.
\end{enumerate}

\begin{fmnote}\fmuline{NOTA}: perquè els vectors \fmm{\vec v_1,\vec v_2,\cdots,\vec v_n} siguen linealment independents, l'expressió \fmm{\lambda_1\cdot\vec v_1+\lambda_2\cdot\vec v_2+\cdots+\lambda_n\cdot\vec v_n=0} ha de tindre una única solució \fmm{\lambda_1=\lambda_2=\cdots=\lambda_n=0}.\end{fmnote}
\end{fmdef}

\Needspace{7\baselineskip}
\begin{fmex}{Exemples}
\begin{enumerate}[label=\textbf{\arabic*.},leftmargin=*]
\item \textbf{Els vectors \fmm{(1,1,4),(1,2,3)} i \fmm{(0,3,7)} són linealment independents (l.i.)?}

\fmdisp{\lambda_1(1,1,4)+\lambda_2(1,2,3)+\lambda_3(0,3,7)=(0,0,0)}

\fmdisp{\left.\begin{aligned}\lambda_1+\lambda_2&=0\\ \lambda_1+2\lambda_2+3\lambda_3&=0\\ 4\lambda_1+3\lambda_2+7\lambda_3&=0\end{aligned}\right\}\longrightarrow\left\{\begin{aligned}\lambda_1&=-\lambda_2\\ -\lambda_2+2\lambda_2+3\lambda_3&=0\\ -4\lambda_2+3\lambda_2+7\lambda_3&=0\end{aligned}\right.\longrightarrow\left\{\begin{aligned}\lambda_2+3\lambda_3&=0\\ -\lambda_2+7\lambda_3&=0\end{aligned}\right.}

\fmdisp{\xrightarrow[\text{sumem equacions}]{}\ 10\lambda_3=0\rightarrow\boxed{\lambda_3=0}\rightarrow\lambda_2+3\cdot0=0\rightarrow\boxed{\lambda_2=0}\rightarrow\boxed{\lambda_1=0}}

Aleshores podem afirmar que els tres vectors són l.i.
\item \textbf{Els vectors \fmm{(1,1,4),(1,2,3)} i \fmm{(2,0,10)} són linealment independents (l.i.)?}

\fmdisp{\lambda_1(1,1,4)+\lambda_2(1,2,3)+\lambda_3(2,0,10)=(0,0,0)}

\fmdisp{\left.\begin{aligned}\lambda_1+\lambda_2+2\lambda_3&=0\\ \lambda_1+2\lambda_2&=0\\ 4\lambda_1+3\lambda_2+10\lambda_3&=0\end{aligned}\right\}\rightarrow\left\{\begin{aligned}-2\lambda_2+\lambda_2+2\lambda_3&=0\\ \lambda_1&=-2\lambda_2\\ -8\lambda_2+3\lambda_2+10\lambda_3&=0\end{aligned}\right.\rightarrow\left\{\begin{aligned}-\lambda_2+2\lambda_3&=0\\ -5\lambda_2+10\lambda_3&=0\end{aligned}\right.}

Si resolem aquest sistema d'equacions veiem que és un SCI (la segona equació és la primera multiplicada per 5); com té més d'una solució, podem afirmar que els vectors són l.d.
\end{enumerate}
\end{fmex}

Seguidament aplicarem aquests conceptes de vectors als determinants i les matrius, per tal d'introduir el rang d'una matriu. Aquest concepte veurem que està molt relacionat amb la dependència lineal entre els vectors que formen una matriu.

\Needspace{7\baselineskip}
\begin{fmprop}{Dependència o independència amb determinants}
Amb tot el que hem vist abans podem afirmar que si un conjunt de vectors són l.d., podem expressar un d'ells com a combinació lineal dels altres: \fmm{\vec v_1=\lambda_2\cdot\vec v_2+\lambda_3\cdot\vec v_3+\cdots+\lambda_n\cdot\vec v_n}.\newline Si pensem en un determinant format per vectors fila tenim:

\fmdisp{\begin{vmatrix}\vec v_1\\ \vec v_2\\ \vdots\\ \vec v_n\end{vmatrix}=\begin{vmatrix}\lambda_2\cdot\vec v_2+\cdots+\lambda_n\cdot\vec v_n\\ \vec v_2\\ \vdots\\ \vec v_n\end{vmatrix}=\lambda_2\underbrace{\begin{vmatrix}\vec v_2\\ \vec v_2\\ \vdots\\ \vec v_n\end{vmatrix}}_{0}+\cdots+\lambda_n\underbrace{\begin{vmatrix}\vec v_n\\ \vec v_2\\ \vdots\\ \vec v_n\end{vmatrix}}_{0}=0}

ja que si dues files són iguals el determinant val zero.\newline Aleshores podem afirmar que:

\begin{itemize}
\item \textbf{Si la matriu \fmm{A} és regular, \fmm{|A|\neq0}, tots els seus vectors fila o columna seran independents.}
\item \textbf{Si la matriu \fmm{A} és singular, \fmm{|A|=0}, tots els seus vectors fila o columna seran dependents.}
\end{itemize}

D'aquesta forma, per saber si un conjunt de vectors són l.i. o l.d. calcularem el determinant de la matriu formada pels vectors, que és més senzill que resoldre un sistema.

Però el problema apareix si la matriu formada pels vectors no és una matriu quadrada; per exemple, si volem saber si els vectors següents són l.i. o l.d.:\newline -- (2,1), (4,5) i (1,3).\newline -- (2,1,3,1), (4,1,2,3) i (2,1,0,1).

Per a resoldre aquest problema utilitzem els rangs de les matrius.
\end{fmprop}

\Needspace{7\baselineskip}
\begin{fmdef}{Rang d'una matriu}
El \textbf{rang d'una matriu} \fmm{A} de qualsevol ordre és el nombre màxim de vectors fila, o de vectors columna, d'\fmm{A} que són linealment independents. Ho escriurem com \textbf{rg(A)}.

\begin{fmnote}El rang d'una matriu \fmm{A} coincideix amb l'ordre de la major submatriu quadrada el determinant de la qual és diferent de zero.\end{fmnote}
\end{fmdef}

\Needspace{7\baselineskip}
\begin{fmprop}{Calculem el rang d'una matriu}
\begin{enumerate}
\item Suprimir les línies dependents, és a dir, nu\lgem{}es, repetides, múltiples o combinacions lineals evidents.
\item Trobar el determinant de la matriu amb major dimensió possible en la matriu \fmm{A}. Si és distint de zero direm que: \textbf{rg(A) = dimensió del determinant}.
\item Si el determinant és zero, se n'ha de calcular, si és possible, un altre amb la mateixa dimensió. Si s'esgoten totes les possibilitats d'una determinada dimensió, es calculen els determinants de les submatrius d'una dimensió inferior.
\end{enumerate}
\end{fmprop}

\Needspace{7\baselineskip}
\begin{fmex}{Exemples de càlcul del rang de matrius}
\begin{enumerate}[label=\textbf{\arabic*.},leftmargin=*]
\item Calculem el rang de la matriu següent: \fmm{A=\begin{pmatrix}-1 & 3 & 0 & 1 & 2\\ 0 & 5 & 1 & 2 & 3\\ -3 & -1 & -2 & -1 & 0\\ 3 & 11 & 4 & 5 & 6\end{pmatrix}}

Comencem per destacar un menor d'ordre dos no nul (això es pot fer a ull). Procurem que siga tan senzill com siga possible, ja que formarà part de molts determinants d'ordre superior que haurem de calcular més avant.

\fmdisp{A=\begin{pmatrix}\boxed{-1} & 3 & \boxed{0} & 1 & 2\\ \boxed{0} & 5 & \boxed{1} & 2 & 3\\ -3 & -1 & -2 & -1 & 0\\ 3 & 11 & 4 & 5 & 6\end{pmatrix}\qquad\begin{vmatrix}-1 & 0\\ 0 & 1\end{vmatrix}\neq0}

El fet que aquest menor siga no nul ens assegura que les dues primeres files d'\fmm{A} són l.i.

Vegem si la tercera fila depèn linealment d'aquestes. Per a això, hi afegim els dos elements nous, \fmm{-3} i \fmm{-2}, i calculem tots els menors d'ordre 3 en què intervinguen aquestes dues columnes, per si algun és diferent de zero:

\fmdisp{\begin{vmatrix}\color[HTML]{e5484d}-1 & 3 & \color[HTML]{e5484d}0\\ \color[HTML]{e5484d}0 & 5 & \color[HTML]{e5484d}1\\ -3 & -1 & -2\end{vmatrix}=0\qquad\begin{vmatrix}\color[HTML]{e5484d}-1 & \color[HTML]{e5484d}0 & 1\\ \color[HTML]{e5484d}0 & \color[HTML]{e5484d}1 & 2\\ -3 & -2 & -1\end{vmatrix}=0\qquad\begin{vmatrix}\color[HTML]{e5484d}-1 & \color[HTML]{e5484d}0 & 2\\ \color[HTML]{e5484d}0 & \color[HTML]{e5484d}1 & 3\\ -3 & -2 & 0\end{vmatrix}=0}

En ser zero tots aquests, es dedueix que \textbf{la tercera fila és combinació lineal de les dues primeres.}

(Si algun d'aquests tres menors fora no nul, la tercera fila seria l.i. de les dues primeres. Prendríem com a base aquest menor d'ordre 3 i ampliaríem les tres columnes amb els elements corresponents de la quarta fila. Obtindríem així tres columnes de quatre elements. Afegint-hi una altra columna, provaríem si algun dels menors d'ordre 4 era no nul.)

Ara farem amb la quarta fila el mateix que hem fet amb la tercera. Els tres menors d'ordre 3 que així s'obtenen són zero. Per tant, \textbf{la quarta fila és, també, combinació lineal de les dues primeres.}

El rang de la matriu és dos (màxim ordre d'un menor no nul o, el que és el mateix, màxim nombre de files l.i.). Aleshores \fmm{\boxed{rg(A)=2}}.
\item Siga la matriu \fmm{A=\begin{pmatrix}1 & 1 & 2\\ 1 & 2 & 0\\ 4 & 3 & 10\end{pmatrix}}, calculem el rang:

Calculem el determinant: \fmm{|A|=\begin{vmatrix}1 & 1 & 2\\ 1 & 2 & 0\\ 4 & 3 & 10\end{vmatrix}=20+6+0-(16+0+10)=0}, aleshores el rang no pot ser 3. Busquem un menor d'ordre 2x2 diferent de zero:

\fmdisp{\begin{vmatrix}1 & 1\\ 1 & 2\end{vmatrix}=2-1=1\neq0\rightarrow\boxed{rg(A)=2}}

Aleshores (1,1,2) i (1,2,0) són l.i., però (1,1,2), (1,2,0) i (4,3,10) són l.d.
\item Estudiem el rang de la matriu: \fmm{B=\begin{pmatrix}1 & 1 & 0\\ 1 & 2 & 3\\ 4 & 3 & 7\end{pmatrix}}

\fmdisp{|B|=\begin{vmatrix}1 & 1 & 0\\ 1 & 2 & 3\\ 4 & 3 & 7\end{vmatrix}=14+0+12-(0+9+7)=10\neq0\rightarrow\boxed{rg(B)=3}}

\fmm{\rightarrow} els tres vectors són l.i.
\item Calculem el rang de la matriu: \fmm{A=\begin{pmatrix}1 & 3 & 1 & 2 & 4 & 0\\ 2 & 1 & -3 & 4 & 3 & 0\\ 3 & 0 & -6 & 6 & 3 & 0\\ 1 & 2 & 0 & 2 & 3 & 0\end{pmatrix}}

Si ens fixem, veiem que \fmm{C_4=2\cdot C_1}, \fmm{C_5=C_1+C_2} i \fmm{C_6} és un vector nul; aleshores podem eliminar aquestes columnes per a estudiar el rang de la matriu \fmm{A}, ja que sabem que són l.d. de les altres.

\fmdisp{rg(A)=rg\begin{pmatrix}1 & 3 & 1\\ 2 & 1 & -3\\ 3 & 0 & -6\\ 1 & 2 & 0\end{pmatrix}}

Aleshores agafem menors d'ordre 3x3 per veure si n'hi ha algun diferent de zero:

\fmdisp{\begin{vmatrix}F_1\\ F_2\\ F_3\end{vmatrix}=\begin{vmatrix}1 & 3 & 1\\ 2 & 1 & -3\\ 3 & 0 & -6\end{vmatrix}=(-6+0-27)-(3-36)=-33+33=0}

\fmdisp{\begin{vmatrix}F_1\\ F_2\\ F_4\end{vmatrix}=\begin{vmatrix}1 & 3 & 1\\ 2 & 1 & -3\\ 1 & 2 & 0\end{vmatrix}=(0+4-9)-(1-6+0)=-5+5=0}

\fmdisp{\begin{vmatrix}F_2\\ F_3\\ F_4\end{vmatrix}=\begin{vmatrix}2 & 1 & -3\\ 3 & 0 & -6\\ 1 & 2 & 0\end{vmatrix}=(0-18-6)-(-24)=-24+24=0}

\fmdisp{\begin{vmatrix}F_1\\ F_3\\ F_4\end{vmatrix}=\begin{vmatrix}1 & 3 & 1\\ 3 & 0 & -6\\ 1 & 2 & 0\end{vmatrix}=(0+6-18)-(-12)=-12+12=0}

Com tots els menors d'ordre 3 són zero, el rang de la matriu \fmm{A} serà 2 com a màxim. Comprovem-ho:

\fmdisp{\begin{vmatrix}1 & 3\\ 2 & 1\end{vmatrix}=1-6=-5\neq0\longrightarrow\boxed{rg(A)=2}}

\end{enumerate}
\end{fmex}

\Needspace{7\baselineskip}
\begin{fmprop}{Discussió de rangs en funció d'un paràmetre}
Consisteix a esbrinar els diferents rangs que pot tindre una matriu segons els valors que prenen alguns dels seus elements, que venen expressats en forma de paràmetres:

\begin{itemize}
\item Si la matriu és quadrada, es resol l'equació \fmm{|A|=0}. Les solucions són els diferents casos que cal estudiar.
\item Si la matriu és rectangular, es resol l'equació \fmm{|A'|=0}, on \fmm{A'} és una submatriu quadrada d'\fmm{A} de la major grandària possible. Les solucions són els casos que cal estudiar.\newline Tant se val la submatriu que triem: encara que no isquen els mateixos casos, arribarem a les mateixes conclusions.
\end{itemize}
\end{fmprop}

\Needspace{7\baselineskip}
\begin{fmex}{Exemples d'estudi del rang de matrius amb paràmetres}
\begin{enumerate}[label=\textbf{\arabic*.},leftmargin=*]
\item Siga la matriu \fmm{A=\begin{pmatrix}a & 2\\ 2 & a\end{pmatrix}\rightarrow|A|=a^2-4=0\rightarrow a=\pm2}\newline
          Així, si \fmm{\boxed{a=2\rightarrow rg(A)=1}}, si \fmm{\boxed{a=-2\rightarrow rg(A)=1}} i si \fmm{\boxed{a\neq\pm2\rightarrow rg(A)=2}}
\item Siga la matriu \fmm{B=\begin{pmatrix}a & 2 & 2\\ 2 & a & 2\\ 2 & 2 & a\end{pmatrix}}; com abans, calculem el determinant:

\fmm{|B|=a^3+8+8-(4a+4a+4a)=a^3-12a+16=0}; si resolem l'equació per Ruffini obtenim les solucions: \fmm{a=2} (arrel doble) i \fmm{a=-4}. Estudiem els diferents casos:

\begin{itemize}
\item Si \fmm{a=2\rightarrow\begin{pmatrix}2 & 2 & 2\\ 2 & 2 & 2\\ 2 & 2 & 2\end{pmatrix}\rightarrow\boxed{rg(B)=1}}
\item Si \fmm{a=-4\rightarrow|B|=0} estudiem un menor d'ordre 2: \fmm{\begin{vmatrix}-4 & 2\\ 2 & -4\end{vmatrix}=16-4=12\neq0\rightarrow\boxed{rg(B)=2}}
\item Si \fmm{a\neq2} i \fmm{a\neq-4\rightarrow\boxed{rg(B)=3}}
\end{itemize}
\item Siga la matriu \fmm{C=\begin{pmatrix}a & 2 & 4\\ 1 & 1 & a\end{pmatrix}}. Per estudiar el rang, agafem els diferents menors d'ordre 2:

\fmdisp{C_1'=\begin{pmatrix}a & 2\\ 1 & 1\end{pmatrix}\rightarrow|C_1'|=a-2=0\rightarrow a=2}

\fmdisp{C_2'=\begin{pmatrix}2 & 4\\ 1 & a\end{pmatrix}\rightarrow|C_2'|=2a-4=0\rightarrow a=2}

\fmdisp{C_3'=\begin{pmatrix}a & 4\\ 1 & a\end{pmatrix}\rightarrow|C_3'|=a^2-4=0\rightarrow a=\pm2}

\newpage
Aleshores:

\begin{itemize}
\item Si \fmm{\boxed{a=2\rightarrow rg(C)=1}}
\item Si \fmm{\boxed{a\neq2\rightarrow rg(C)=2}}
\end{itemize}
\end{enumerate}
\end{fmex}

\Needspace{18\baselineskip}
\section{Matriu inversa}\hypertarget{sec6}{}

\Needspace{7\baselineskip}
\begin{fmdef}{Matriu inversa}
Siga \fmm{A} una matriu quadrada regular d'ordre \fmm{n} (\fmm{|A_n|\neq0}); s'anomena \textbf{matriu inversa d'\fmm{A}} una altra matriu, \fmm{A^{-1}}, del mateix ordre, que verifica:

\fmdisp{\boxed{A\cdot A^{-1}=A^{-1}\cdot A=I_n}}

Per tant, una matriu \fmuline{té inversa o és invertible si el seu determinant és diferent de zero}, és a dir, si és regular.

A més, tenim una fórmula per poder calcular la inversa d'una matriu \fmm{A}:

\fmdisp{\boxed{A^{-1}=\frac{1}{|A|}\cdot\left(Adj(A)\right)^t}}

\end{fmdef}

\Needspace{7\baselineskip}
\begin{fmex}{Formes de calcular la matriu inversa d'A}
\begin{enumerate}[label=\textbf{\arabic*.},leftmargin=*]
\item \textbf{Utilitzant la definició:} \fmm{A\cdot A^{-1}=A^{-1}\cdot A=I_n}\newline
          Vegem un exemple: donada la matriu \fmm{A=\begin{pmatrix}1 & 2\\ 3 & 5\end{pmatrix}}, calculem la seua inversa \fmm{A^{-1}=\begin{pmatrix}a & b\\ c & d\end{pmatrix}}.

\fmdisp{A\cdot A^{-1}=\begin{pmatrix}1 & 2\\ 3 & 5\end{pmatrix}\cdot\begin{pmatrix}a & b\\ c & d\end{pmatrix}=\begin{pmatrix}a+2c & b+2d\\ 3a+5c & 3b+5d\end{pmatrix}=\begin{pmatrix}1 & 0\\ 0 & 1\end{pmatrix}}

Ara resolem el sistema de 4 equacions amb 4 incògnites:

\fmdisp{\left.\begin{aligned}a+2c&=1\\ 3a+5c&=0\end{aligned}\right\}\rightarrow\left.\begin{aligned}-3a-6c&=-3\\ 3a+5c&=0\end{aligned}\right\}\rightarrow-c=-3\rightarrow\boxed{c=3}\rightarrow\boxed{a=-5}}

\fmdisp{\left.\begin{aligned}b+2d&=0\\ 3b+5d&=1\end{aligned}\right\}\rightarrow\left.\begin{aligned}-3b-6d&=0\\ 3b+5d&=1\end{aligned}\right\}\rightarrow-d=1\rightarrow\boxed{d=-1}\rightarrow\boxed{b=2}}

Aleshores:

\fmdisp{\boxed{A^{-1}=\begin{pmatrix}-5 & 2\\ 3 & -1\end{pmatrix}}}

\item \textbf{A partir de la fórmula:} \fmm{A^{-1}=\dfrac{1}{|A|}\cdot(Adj(A))^t}\newline
          Calculem la matriu inversa de \fmm{A=\begin{pmatrix}3 & 1 & 2\\ 1 & 3 & 1\\ 2 & 4 & 2\end{pmatrix}}

\fmdisp{|A|=\begin{vmatrix}3 & 1 & 2\\ 1 & 3 & 1\\ 2 & 4 & 2\end{vmatrix}=(18+8+2)-(12+12+2)=28-26=2}

Ara calculem la matriu adjunta (Adj(A)):

\fmdisp{A_{11}=\begin{vmatrix}3 & 1\\ 4 & 2\end{vmatrix}\cdot(-1)^{1+1}=2\quad A_{12}=\begin{vmatrix}1 & 1\\ 2 & 2\end{vmatrix}\cdot(-1)^{1+2}=0\quad A_{13}=\begin{vmatrix}1 & 3\\ 2 & 4\end{vmatrix}\cdot(-1)^{1+3}=-2}

\fmdisp{A_{21}=\begin{vmatrix}1 & 2\\ 4 & 2\end{vmatrix}\cdot(-1)^{2+1}=6\quad A_{22}=\begin{vmatrix}3 & 2\\ 2 & 2\end{vmatrix}\cdot(-1)^{2+2}=2\quad A_{23}=\begin{vmatrix}3 & 1\\ 2 & 4\end{vmatrix}\cdot(-1)^{2+3}=-10}

\fmdisp{A_{31}=\begin{vmatrix}1 & 2\\ 3 & 1\end{vmatrix}\cdot(-1)^{3+1}=-5\quad A_{32}=\begin{vmatrix}3 & 2\\ 1 & 1\end{vmatrix}\cdot(-1)^{3+2}=-1\quad A_{33}=\begin{vmatrix}3 & 1\\ 1 & 3\end{vmatrix}\cdot(-1)^{3+3}=8}

Aleshores tenim que:

\fmdisp{(Adj(A))^t=\begin{pmatrix}2 & 0 & -2\\ 6 & 2 & -10\\ -5 & -1 & 8\end{pmatrix}^t=\begin{pmatrix}2 & 6 & -5\\ 0 & 2 & -1\\ -2 & -10 & 8\end{pmatrix}}

Aplicant la fórmula:

\fmdisp{\boxed{A^{-1}=\frac12\begin{pmatrix}2 & 6 & -5\\ 0 & 2 & -1\\ -2 & -10 & 8\end{pmatrix}=\begin{pmatrix}1 & 3 & -5/2\\ 0 & 1 & -1/2\\ -1 & -5 & 4\end{pmatrix}}}

\item \textbf{Mitjançant el mètode de Gauss.}\newline
          Co\lgem{}oquem la matriu \fmm{A} i a la dreta la matriu \fmm{I}, formant una matriu \fmm{n\times2n}. Realitzem les transformacions necessàries per a obtindre la matriu \fmm{I} en el lloc on estava \fmm{A}, i la matriu que ens queda en el lloc d'\fmm{I} és la matriu \fmm{A^{-1}}.\newline
          Vegem un exemple. Calculem la matriu inversa de \fmm{A=\begin{pmatrix}3 & 1 & 2\\ 1 & 3 & 1\\ 2 & 4 & 2\end{pmatrix}}

\fmdisp{\left(\begin{array}{ccc|ccc}3 & 1 & 2 & 1 & 0 & 0\\ 1 & 3 & 1 & 0 & 1 & 0\\ 2 & 4 & 2 & 0 & 0 & 1\end{array}\right)\xrightarrow[F_1\leftrightarrow F_2]{}\left(\begin{array}{ccc|ccc}1 & 3 & 1 & 0 & 1 & 0\\ 3 & 1 & 2 & 1 & 0 & 0\\ 2 & 4 & 2 & 0 & 0 & 1\end{array}\right)}

\fmdisp{\xrightarrow[F_2-3F_1]{F_3-2F_1}\left(\begin{array}{ccc|ccc}1 & 3 & 1 & 0 & 1 & 0\\ 0 & -8 & -1 & 1 & -3 & 0\\ 0 & -2 & 0 & 0 & -2 & 1\end{array}\right)\xrightarrow[F_2\leftrightarrow F_3]{}\left(\begin{array}{ccc|ccc}1 & 3 & 1 & 0 & 1 & 0\\ 0 & -2 & 0 & 0 & -2 & 1\\ 0 & -8 & -1 & 1 & -3 & 0\end{array}\right)}

\fmdisp{\xrightarrow[F_3-4F_2]{}\left(\begin{array}{ccc|ccc}1 & 3 & 1 & 0 & 1 & 0\\ 0 & -2 & 0 & 0 & -2 & 1\\ 0 & 0 & -1 & 1 & 5 & -4\end{array}\right)\xrightarrow[F_2=(-1/2)F_2]{F_3=-F_3}\left(\begin{array}{ccc|ccc}1 & 3 & 1 & 0 & 1 & 0\\ 0 & 1 & 0 & 0 & 1 & -1/2\\ 0 & 0 & 1 & -1 & -5 & 4\end{array}\right)}

\fmdisp{\xrightarrow[F_1=F_1-3F_2-F_3]{}\left(\begin{array}{ccc|ccc}1 & 0 & 0 & 1 & 3 & -5/2\\ 0 & 1 & 0 & 0 & 1 & -1/2\\ 0 & 0 & 1 & -1 & -5 & 4\end{array}\right)}

Aleshores:

\fmdisp{\boxed{A^{-1}=\begin{pmatrix}1 & 3 & -5/2\\ 0 & 1 & -1/2\\ -1 & -5 & 4\end{pmatrix}}}

\end{enumerate}
\end{fmex}

\Needspace{7\baselineskip}
\begin{fmprop}{Propietats de la matriu inversa}
\begin{multicols}{2}
\begin{itemize}
\item \fmm{(A\cdot B)^{-1}=B^{-1}\cdot A^{-1}}
\item \fmm{(A^{-1})^{-1}=A}
\item \fmm{(A^{-1})^t=(A^t)^{-1}}
\item \fmm{(A^2)^{-1}=(A^{-1})^2}
\item \fmm{(kA)^{-1} = k^{-1}A^{-1}=\frac{1}{k}A^{-1}}
\item \fmm{|A^{-1}|=\dfrac{1}{|A|}}
\end{itemize}
\end{multicols}
\end{fmprop}

\Needspace{18\baselineskip}
\section{Equacions i sistemes matricials}\hypertarget{sec7}{}

\Needspace{7\baselineskip}
\begin{fmprop}{Resolem equacions i sistemes matricials}
Són equacions o sistemes en els quals apareixen termes que són matrius. La idea és que els resoldrem igual que les equacions i sistemes amb nombres reals. Però \fmuline{cal recordar} que amb les matrius tenim certes restriccions:

\begin{enumerate}
\item No existeix la divisió de matrius. Per tant, utilitzarem la matriu inversa (sempre que existisca).
\item No es compleix la propietat commutativa per al producte de matrius.
\end{enumerate}
\end{fmprop}

\Needspace{7\baselineskip}
\begin{fmex}{Exemples d'equacions i sistemes}
\begin{enumerate}[label=\textbf{\alph*)},leftmargin=*]
\item Donades les matrius \fmm{A=\begin{pmatrix}1 & 2\\ 3 & 5\end{pmatrix}}, \fmm{B=\begin{pmatrix}3 & 4\\ 5 & 1\end{pmatrix}} i \fmm{C=\begin{pmatrix}0 & 2\\ 4 & 6\end{pmatrix}}, resolem \fmm{\boxed{AX-2B=3C}}:

\fmdisp{AX=3C+2B\rightarrow A^{-1}\cdot A\cdot X=A^{-1}\cdot(3C+2B)\rightarrow X=A^{-1}(3C+2B)}

\fmdisp{3C+2B=3\cdot\begin{pmatrix}0 & 2\\ 4 & 6\end{pmatrix}+2\cdot\begin{pmatrix}3 & 4\\ 5 & 1\end{pmatrix}=\begin{pmatrix}6 & 14\\ 22 & 20\end{pmatrix}}

Ara calcularem \fmm{A^{-1}}: primer \fmm{|A|=5-6=-1\rightarrow A} és invertible.

\fmdisp{A^{-1}=\frac{1}{-1}\cdot\begin{pmatrix}5 & -3\\ -2 & 1\end{pmatrix}^t=\begin{pmatrix}-5 & 2\\ 3 & -1\end{pmatrix}}

\fmdisp{X=A^{-1}(3C+2B)=\begin{pmatrix}-5 & 2\\ 3 & -1\end{pmatrix}\cdot\begin{pmatrix}6 & 14\\ 22 & 20\end{pmatrix}=\begin{pmatrix}-30+44 & -70+40\\ 18-22 & 42-20\end{pmatrix}}

\fmdisp{\boxed{X=\begin{pmatrix}14 & -30\\ -4 & 22\end{pmatrix}}}

\item Donades les matrius \fmm{A=\begin{pmatrix}4 & 2\\ 3 & 5\end{pmatrix}} i \fmm{B=\begin{pmatrix}3 & 4\\ 5 & 1\end{pmatrix}}, resolem el sistema: \fmm{\left\{\begin{aligned}2X+Y&=A\\ 4X-3Y&=B\end{aligned}\right.}

\fmdisp{\left.\begin{aligned}2X+Y&=A\\ 4X-3Y&=B\end{aligned}\right\}\xrightarrow[Y=A-2X]{}4X-3(A-2X)=B\rightarrow4X-3A+6X=B}

\fmdisp{10X=B+3A\rightarrow X=\frac{B+3A}{10}}

\fmdisp{Y=A-2\left(\frac{B+3A}{10}\right)=\frac{5A}{5}-\frac{B}{5}-\frac{3A}{5}=\frac{2A-B}{5}}

Fem les operacions de matrius per a calcular \fmm{X} i \fmm{Y} i obtenim que:

\fmdisp{\boxed{X=\begin{pmatrix}3/2 & 1\\ 7/5 & 8/5\end{pmatrix};\ Y=\begin{pmatrix}1 & 0\\ 1/5 & 9/5\end{pmatrix}}}

\item Donada la matriu \fmm{A=\begin{pmatrix}3 & 0\\ 1 & 3\end{pmatrix}}, calculem la matriu \fmm{X} que verifica aquesta equació:

\fmdisp{XA-2X=A}

\newpage
Primer pas: aïllem \fmm{X} traient factor comú en el primer membre, tenint en compte que:

\fmdisp{2X=2XI=2IX\rightarrow X(A-2I)=A}

Multipliquem els dos membres per \fmm{(A-2I)^{-1}} per la dreta.\newline Per això, calculem \fmm{(A-2I)=\begin{pmatrix}1 & 0\\ 1 & 1\end{pmatrix}} i la seua inversa, \fmm{(A-2I)^{-1}=\begin{pmatrix}1 & 0\\ 1 & 1\end{pmatrix}^{-1}=\begin{pmatrix}1 & 0\\ -1 & 1\end{pmatrix}}

\fmdisp{X=A\cdot(A-2I)^{-1}\rightarrow X=\begin{pmatrix}3 & 0\\ 1 & 3\end{pmatrix}\cdot\begin{pmatrix}1 & 0\\ -1 & 1\end{pmatrix}=\boxed{\begin{pmatrix}3 & 0\\ -2 & 3\end{pmatrix}}}

\end{enumerate}
\end{fmex}

\end{document}
